\documentclass[a4paper,11pt]{article}
\usepackage[T1]{fontenc}
\usepackage[utf8]{inputenc}
\usepackage[ngerman]{babel}
\usepackage[margin=2.5cm]{geometry}

\begin{document}

\section*{Blatt 3, Aufgabe 3.1.4: Konfidenzintervalle}

\textbf{Bedeutung eines zweiseitigen Konfidenzintervalls für
\(\alpha=0{,}05\).}
Ein zweiseitiges Konfidenzintervall hat das Konfidenzniveau
\(1-\alpha=0{,}95\), also \(95\,\%\). Die Irrtumswahrscheinlichkeit
\(\alpha\) wird gleichmäßig auf beide Seiten verteilt:
\(\alpha/2=0{,}025\) liegt jeweils im linken und im rechten Rand der
Verteilung. Bei wiederholter Anwendung des Verfahrens enthalten daher
\(95\,\%\) der so berechneten Intervalle den wahren Parameter. Für den
unbekannten Mittelwert lautet das Intervall bei Verwendung der
Student-\(t\)-Verteilung beispielsweise
\[
\bar{x}\mathbin{\pm}t_{n-1,\,1-\alpha/2}\frac{s}{\sqrt{n}}.
\]

\medskip
\noindent Für \(\alpha=0{,}05\) gilt \(1-\alpha/2=0{,}975\). Die
benötigten Quantile der Student-\(t\)-Verteilung sind:
\[
\begin{array}{c|c|c}
n & \text{Freiheitsgrade }n-1 & t_{n-1,\,0{,}975}\\
\hline
2 & 1 & 12{,}706\\
3 & 2 & 4{,}303\\
11 & 10 & 2{,}228\\
101 & 100 & 1{,}984
\end{array}
\]

\end{document}
